\documentclass[aps,pra,twocolumn,superscriptaddress,10pt]{revtex4-2}

\usepackage{amsmath}
\usepackage{amssymb}
\usepackage{graphicx}
\usepackage{amsthm}
\usepackage{mathtools}

\newtheorem{theorem}{Theorem}
\newtheorem{lemma}[theorem]{Lemma}
\newtheorem{proposition}[theorem]{Proposition}
\newtheorem{definition}[theorem]{Definition}
\newtheorem{remark}[theorem]{Remark}

\DeclareMathOperator{\spanop}{span}
\newcommand{\C}{\mathbb{C}}
\newcommand{\R}{\mathbb{R}}
\newcommand{\PP}{\mathbb{P}}
\newcommand{\YO}{\mathrm{YO}}

\usepackage[
  breaklinks=true,
  colorlinks=true,
  citecolor=blue,
  linkcolor=blue,
  urlcolor=blue
]{hyperref}
\usepackage{orcidlink}

\begin{document}

\title{Chromatic Completeness Does Not Determine Geometric Coordinatizability}

\author{Karl Svozil\,\orcidlink{0000-0001-6554-2802}}
\affiliation{Institute for Theoretical Physics, TU Wien,
Wiedner Hauptstra{\ss}e 8-10/136, A-1040 Vienna, Austria}
\email{karl.svozil@tuwien.ac.at}

\date{\today}

\begin{abstract}
We distinguish two properties of context hypergraphs: \emph{chromatic
completeness}, the existence of a single globally consistent nondegenerate
spectral labeling, and \emph{geometric coordinatizability}, the existence of
an injective orthogonal realization by rays.  A strong chromatic number larger
than the Hilbert-space dimension precludes only the former.  It does not, by
itself, preclude an orthogonal realization.  This is shown by two
three-dimensional hypergraphs with the same strong chromatic number.  A
completed 25-ray Yu--Oh hypergraph has strong chromatic number four and an
explicit edge-faithful realization in $\R^3$.  Greechie's $G_{32}$ also has
strong chromatic number four, as well as a separating and unital set of
two-valued states, but admits no injective orthogonal realization in $\C^3$.
We give an elementary algebraic proof: its incidence relations force two
distinct vertices to represent the same ray.  Thus the strong chromatic number
does not determine geometric coordinatizability.
\end{abstract}

\maketitle

\section{Introduction}

An $n$-uniform \emph{context hypergraph} $H=(V,E)$ consists of a vertex set
$V$ and a collection $E$ of $n$-element hyperedges, called \emph{contexts}.
It is an abstract incidence structure until a ray realization is specified.
In the intended three-dimensional case a context represents a triple of
mutually orthogonal rank-one projectors, or equivalently a triple of mutually
orthogonal rays.

A \emph{faithful orthogonal representation} (FOR) of an $n$-uniform
hypergraph $H$ in $\C^n$ is an injective assignment
\[
    i\longmapsto [v_i]\in \PP(\C^n)
\]
such that the rays assigned to every context are pairwise orthogonal, and
hence admit normalized representatives forming an orthonormal basis.  Thus
\[
    \rho(i)\perp\rho(j)\qquad
    \text{whenever }\{i,j\}\subseteq e\in E.
\]
Throughout, \emph{faithful} means injective on rays; it does not additionally
require nonincident vertex pairs to be nonorthogonal.  When the latter
stronger property holds, we say that the realization is \emph{edge-faithful}.
The main negative result below uses only injectivity.  The existence of
orthogonal ray realizations is central to the quantum logic of the
Kochen--Specker theorem~\cite{kochen1}.

Recent work on chromatic contextuality~\cite{svozil-2021-chroma,svozil-2025-color} emphasizes that if a hypergraph requires more than $n$
colors, then it cannot support a representation in which every context is
assigned the same fixed set of $n$ spectral labels.  This observation is
correct, but it concerns an additional layer of structure: the global
spectral labeling of an already specified incidence geometry.  It does not
settle whether the incidence structure itself can be coordinatized by rays.

The purpose of this paper is to separate these two issues.  Chromatic
completeness is a global labeling condition.  Geometric coordinatizability is
a projective-geometric realization condition.  These two conditions are
logically independent.  The completed Yu--Oh configuration~\cite{Yu-2012} gives a
coordinatizable hypergraph with strong chromatic number four in dimension
three~\cite{svozil-2025-color}.
Greechie's $G_{32}$~\cite{greechie:71} gives a hypergraph with the same strong chromatic
number~\cite{svozil-2021-chroma}, and even a separating and unital set of two-valued states, but no
faithful representation in $\C^3$.  Thus the chromatic number, the supply of
two-valued states, and geometric coordinatizability are three distinct layers
of structure.

\section{Definitions}

The term ``chromatic number'' for hypergraphs is not completely uniform in
the literature.  Throughout this paper we use the following convention.

\begin{definition}[Strong chromatic number]
Let $H=(V,E)$ be a hypergraph.  A \emph{strong $k$-coloring} is a map
$c:V\to\{1,\ldots,k\}$ such that any two vertices occurring in a common
context have different colors.  Equivalently, for every context
$e\in E$, the restriction $c|_e$ is injective.  The least such $k$ is denoted
by $\chi(H)$.
\end{definition}

For an $n$-uniform hypergraph, a strong $n$-coloring assigns all $n$ colors
to every context exactly once.  This is the coloring notion relevant to
nondegenerate spectral assignments.

\begin{definition}[Chromatic completeness]\label{def:chromatic-completeness}
An $n$-uniform context hypergraph $H=(V,E)$ is \emph{chromatically complete}
if there exist a single set
\[
    \Sigma=\{\lambda_1,\ldots,\lambda_n\}\subset\R
\]
of $n$ distinct scalars and a map $c:V\to\Sigma$ such that for every context
$e\in E$, the restriction $c|_e:e\to\Sigma$ is a bijection.
\end{definition}

Thus each context receives the same spectrum, with each spectral value used
exactly once.

\begin{remark}[Chromatic completeness and strong coloring]\label{rem:chromatic-completeness}
For an $n$-uniform hypergraph, chromatic completeness is exactly strong
$n$-colorability.  The $n$ spectral labels in Definition~\ref{def:chromatic-completeness}
define a strong $n$-coloring, because every context receives each label
exactly once.  Conversely, any strong $n$-coloring may be relabeled by an
arbitrary fixed nondegenerate real spectrum
$\Sigma=\{\lambda_1,\ldots,\lambda_n\}$, giving precisely the spectral
assignment required in Definition~\ref{def:chromatic-completeness}.
\end{remark}

Consequently, if $\chi(H)>n$, then $H$ has no chromatically complete spectral
labeling.  At the level of a fixed projective measurement, the numerical
eigenvalues of a nondegenerate self-adjoint observable are labels of its
rank-one spectral projections: any injective map from its finite spectrum to
$\R$ gives another observable with the same projections.  Chromatic
completeness asks whether these labels can be chosen globally and
context-independently.  Its failure need not signal a failure of the
underlying ray geometry.

\section{Independence of coloring and coordinatization}

The following proposition makes the separation explicit.  The examples are
already available in dimension three.

\begin{proposition}[Independence]\label{prop:independence}
Already for $n=3$, the strong chromatic number does not determine the
existence of a faithful orthogonal representation.  In particular, there are
two 3-uniform hypergraphs with $\chi=4$, one of which has an edge-faithful
representation in $\R^3$ and one of which has no faithful representation in
$\C^3$.
\end{proposition}

\begin{proof}
The completed Yu--Oh hypergraph of Sec.~\ref{sec:yuo} has $\chi=4$ and an
explicit edge-faithful representation in $\R^3$.  Greechie's $G_{32}$ has
$\chi=4$, as shown in Sec.~\ref{sec:g32-chromatic}, whereas
Theorem~\ref{thm:g32} proves that it has no faithful orthogonal representation
in $\C^3$.
\end{proof}

\section{A completed 25-ray Yu--Oh hypergraph}\label{sec:yuo}

The Yu--Oh construction starts from 13 rays in $\R^3$~\cite{Yu-2012} (for extensions see Refs.~\cite{harding2025remarksIJTP,Cabello2025,svozil-2025-MYOCHS}).  In
one standard coordinatization these rays are
\begin{align*}
 z_1&=(1,0,0),& z_2&=(0,1,0),& z_3&=(0,0,1),\\
 y_1^-&=(0,1,-1),& y_2^-&=(1,0,-1),& y_3^-&=(1,-1,0),\\
 y_1^+&=(0,1,1),& y_2^+&=(1,0,1),& y_3^+&=(1,1,0),\\
 h_0&=(1,1,1),& h_1&=(-1,1,1),& h_2&=(1,-1,1), \\
 h_3&=(1,1,-1).
\end{align*}
The partial context hypergraph on these 13 rays admits two-valued states.  For
example, assign value $1$ to vertices $3$, $6$, and $8$, and value $0$ to all
other vertices.  Each of the four complete triads then contains exactly one
vertex of value $1$, and no orthogonal pair has both vertices assigned value
$1$.  The Yu--Oh inequality nevertheless gives a state-independent conflict
with noncontextual models reproducing the relevant quantum predictions~\cite{Yu-2012}.
This distinction is useful here: the existence of such partial value
assignments is separate from both chromatic completeness and ray
coordinatizability.

We use a 3-uniform completion of the Yu--Oh orthogonality graph: for each
orthogonal pair among the original 13 rays, a third ray is included so that
the pair belongs to a full orthogonal triad.  The resulting completed
hypergraph has 25 vertices and 16 contexts, and is the completion used in the
chromatic-contextuality analysis of Ref.~\cite{svozil-2025-color}.

Let $H_{\YO}^{+}=(V_{\YO}^{+},E_{\YO}^{+})$, with
$V_{\YO}^{+}=\{1,\ldots,25\}$, have contexts
%\begin{widetext}
\begin{equation}\label{eq:yuo-blocks}
\begin{aligned}
E_{\YO}&^{+}=\{
\{1,2,3\},
\{7,8,9\},
\{4,5,6\},
\{2,8,5\},\\
&\{3,12,25\},
\{9,12,16\},
\{6,12,17\},
\{3,11,14\},\\
&\{7,11,15\},
\{4,11,24\},
\{9,10,23\},
\{1,10,19\},\\
&\{4,10,18\},
\{1,13,20\},
\{7,13,21\},
\{6,13,22\}
\}.
\end{aligned}
\end{equation}
%\end{widetext}
The first four contexts represent the three $\{y_i^+,z_i,y_i^-\}$ triads
and the coordinate triad $\{z_1,z_2,z_3\}$.  The remaining contexts complete
the orthogonal pairs involving $h_0,h_1,h_2,h_3$.

An explicit coordinatization is obtained by assigning to vertex $i$ the ray
spanned by the following unnormalized vector $r_i\in\R^3$:
\begin{widetext}
\begin{align}\label{eq:yuo-coordinates}
\begin{array}{c|rrrrrrrrrrrrr}
 i      &1&2&3&4&5&6&7&8&9&10&11&12&13\\ \hline
 x_i    &0&1&0&1&0&1&1&0&1& 1&-1&1& 1\\
 y_i    &1&0&1&1&0&-1&0&1&0&-1& 1&1& 1\\
 z_i    &1&0&-1&0&1&0&1&0&-1& 1& 1&1&-1
\end{array}
\\[1ex]
\begin{array}{c|rrrrrrrrrrrr}
 i      &14&15&16&17&18&19&20&21&22&23&24&25\\ \hline
 x_i    &2&1&1&1&1&2&2&1&1&1&1&2\\
 y_i    &1&2&-2&1&-1&1&-1&-2&1&2&-1&-1\\
 z_i    &1&-1&1&-2&-2&-1&1&-1&2&1&2&-1
\end{array}
\end{align}
\end{widetext}
That is, $r_i=(x_i,y_i,z_i)^T$, and the represented ray is $[r_i]$.
The first 13 vertices are labeled as follows:
\begin{align*}
 r_1&=y_1^+,& r_2&=z_1,& r_3&=y_1^-,\\
 r_4&=y_3^+,& r_5&=z_3,& r_6&=y_3^-,\\
 r_7&=y_2^+,& r_8&=z_2,& r_9&=y_2^-,\\
 r_{10}&=h_2,& r_{11}&=h_1,& r_{12}&=h_0,& r_{13}&=h_3,
\end{align*}
where equality means equality of rays.  The remaining twelve vertices are
completion rays.  For instance,
\[
    r_{14}\parallel r_3\times r_{11},\qquad
    r_{15}\parallel r_7\times r_{11},\qquad
    r_{23}\parallel r_9\times r_{10},
\]
and similarly for the other added vertices.

Direct evaluation gives $r_i\cdot r_j=0$ precisely for the pairs $\{i,j\}$
contained in a context of $E_{\YO}^{+}$, and no two vectors in
Eq.~\eqref{eq:yuo-coordinates} are proportional.
Thus the coordinates give an edge-faithful orthogonal representation of
$H_{\YO}^{+}$ in $\R^3$.

The strong chromatic number of this completed hypergraph is four.  A strong
3-coloring of $H_{\YO}^{+}$ would restrict to a proper 3-coloring of the
Yu--Oh 13-ray orthogonality graph, because every orthogonal pair among the
original 13 rays appears in at least one completed triad.  Here is a short
self-contained proof that the latter graph is not 3-colorable.  Fix the color
of $h_0$ to be $0$.  Then $y_1^-,y_2^-,y_3^-$ can use only colors $1$ and $2$.
If they all have the same color, the coordinate triad would have to use three
distinct colors while avoiding that common color, which is impossible.
Otherwise, by permuting indices and exchanging colors $1$ and $2$, take
\[
 (c(y_1^-),c(y_2^-),c(y_3^-))=(1,1,2).
\]
The four triads force either
\begin{align*}
 (c(z_1),c(z_2),c(z_3))&=(0,2,1),\\
 (c(y_1^+),c(y_2^+),c(y_3^+))&=(2,0,0),
\end{align*}
or the same two triples with the first two entries interchanged.  In the
first case the three neighbours $y_2^-$, $y_1^+$, and $y_3^+$ of $h_2$ have
all three colors; in the second case the three neighbours $y_1^-$, $y_2^+$,
and $y_3^+$ of $h_1$ do.  Hence three colors do not suffice.

On the other hand, the following strong 4-coloring shows that four colors do
suffice:
\begin{widetext}
\begin{align}\label{eq:yuo-coloring}
\begin{array}{c|rrrrrrrrrrrrr}
 i    &1&2&3&4&5&6&7&8&9&10&11&12&13\\ \hline
 c(i) &0&2&1&0&3&1&0&1&2&1&2&0&2
\end{array}
\\[1ex]
\begin{array}{c|rrrrrrrrrrrr}
 i    &14&15&16&17&18&19&20&21&22&23&24&25\\ \hline
 c(i) &0&1&1&2&2&2&1&1&0&0&1&2
\end{array}
\end{align}
\end{widetext}
Each context in Eq.~\eqref{eq:yuo-blocks} contains three distinct colors.
Therefore
\[
    \chi(H_{\YO}^{+})=4>3,
\]
although $H_{\YO}^{+}$ has an explicit FOR in $\R^3$.  This is the desired
example showing that chromatic obstruction is not geometric
non-coordinatizability.

\section{Greechie's \texorpdfstring{$G_{32}$}{G32} hypergraph}\label{sec:g32}

We now turn to the converse phenomenon: a hypergraph with the same strong
chromatic number as the completed Yu--Oh example, but with no faithful
orthogonal representation in $\C^3$.

\subsection{The incidence structure}

Greechie's diagram $G_{32}$~\cite[Fig.~6]{greechie:71} has vertex set
$\{1,\ldots,15\}$ and ten blocks:
\begin{equation}\label{eq:g32-blocks}
\begin{aligned}
E(G_{32})=\{&\{1,2,3\},
\{3,4,5\},
\{5,6,7\},
\{7,8,9\},\\
&\{9,10,11\},
\{11,12,1\},
\{4,10,13\},\\
&\{6,12,14\},
\{8,2,15\},
\{13,14,15\}\}.
\end{aligned}
\end{equation}
Each block is intended to represent a triad of mutually orthogonal rays.
For a two-valued state $s:V\to\{0,1\}$, write
$\operatorname{supp}(s)=\{i:s(i)=1\}$.  The following six supports give all
two-valued states of $G_{32}$:
\begin{align*}
&\{3,7,10,12,15\}, &&\{3,6,8,11,13\},
&&\{2,5,9,12,13\},\\
&\{2,4,7,11,14\}, &&\{1,5,8,10,14\},
&&\{1,4,6,9,15\}.
\end{align*}
Each support intersects every block in exactly one vertex.  Together they are
unital and separating, as noted previously in Ref.~\cite{svozil-2021-chroma};
nevertheless, the hypergraph is not 3-colorable.  Thus it separates the
existence of two-valued states from chromatic completeness.  The theorem below
shows that it also has a geometric obstruction.

\subsection{The strong chromatic number of
\texorpdfstring{$G_{32}$}{G32}}\label{sec:g32-chromatic}

The hypergraph $G_{32}$ has strong chromatic number four~\cite{svozil-2021-chroma}.  A strong
4-coloring is
\begin{equation}\label{eq:g32-coloring}
\begin{array}{c|rrrrrrrrrrrrrrr}
 i    &1&2&3&4&5&6&7&8&9&10&11&12&13&14&15\\ \hline
 c(i) &0&1&2&0&1&0&2&0&1&2&3&1&1&2&3.
\end{array}
\end{equation}
A direct check of Eq.~\eqref{eq:g32-blocks} shows that each block contains
three distinct colors.

To see that three colors do not suffice, form the dual incidence graph: its
vertices are the ten blocks of $G_{32}$, and each atom is an edge joining the
two blocks in which it occurs.  Since every atom occurs in exactly two blocks,
this dual is cubic.  Numbering the blocks in Eq.~\eqref{eq:g32-blocks} from
left to right, the atom-labelled dual obtained directly from the incidence
list is the Petersen graph.  A strong 3-coloring
of the atoms would be a proper 3-edge-coloring of this cubic graph, because
the three atoms incident with each block must receive different colors.  The
Petersen graph has chromatic index four, not three.  Hence $G_{32}$ has no
strong 3-coloring, and therefore
\[
    \chi(G_{32})=4.
\]
This agrees with the previous chromatic analysis of $G_{32}$ in
Ref.~\cite{svozil-2021-chroma}.

\subsection{Non-coordinatizability theorem}

\begin{theorem}\label{thm:g32}
The hypergraph $G_{32}$ admits no faithful orthogonal representation in
$\C^3$.
\end{theorem}

\begin{proof}
We use the inner-product convention
\[
    \langle x,y\rangle=x^\dagger y,
\]
so the inner product is conjugate-linear in the first argument and linear in
the second.

Assume, for contradiction, that a faithful orthogonal representation exists.
Since $\{13,14,15\}$ is a block, its three rays form an orthonormal basis.
After applying a global unitary transformation, we may fix
\begin{align*}
    v_{13}&=e_1=\begin{pmatrix}1\\0\\0\end{pmatrix},
    &v_{14}&=e_2=\begin{pmatrix}0\\1\\0\end{pmatrix},\\[-1ex]
    v_{15}&=e_3=\begin{pmatrix}0\\0\\1\end{pmatrix}.
\end{align*}

The three blocks meeting this central triangle force three pairs of vectors
into coordinate two-planes.  From $\{4,10,13\}$ we obtain
\[
    v_4=\begin{pmatrix}0\\a\\b\end{pmatrix},\qquad
    v_{10}=\begin{pmatrix}0\\-\overline b\\\overline a\end{pmatrix},
    \qquad |a|^2+|b|^2=1.
\]
From $\{6,12,14\}$ we obtain
\[
    v_6=\begin{pmatrix}c\\0\\d\end{pmatrix},\qquad
    v_{12}=\begin{pmatrix}-\overline d\\0\\\overline c\end{pmatrix},
    \qquad |c|^2+|d|^2=1.
\]
From $\{8,2,15\}$ we obtain
\[
    v_2=\begin{pmatrix}p\\q\\0\end{pmatrix},\qquad
    v_8=\begin{pmatrix}-\overline q\\\overline p\\0\end{pmatrix},
    \qquad |p|^2+|q|^2=1.
\]
These parametrizations are general: in each coordinate two-plane, once one
normalized ray is chosen, the orthogonal ray is unique up to phase, and the
chosen phases above lose no generality at the level of rays.

We shall use the following elementary projection identity.  If
$\{x,y,z\}$ is an orthonormal basis of $\C^3$, $u\perp x$, and $w\perp z$,
then
\begin{equation}\label{eq:projection-identity}
    \langle u,w\rangle
    =\langle u,y\rangle\langle y,w\rangle.
\end{equation}
Indeed,
\[
    w=x\langle x,w\rangle+y\langle y,w\rangle+z\langle z,w\rangle.
\]
Taking the inner product with $u$ eliminates the first term because
$u\perp x$, while the last coefficient vanishes because $w\perp z$.
Only the middle term remains, proving Eq.~\eqref{eq:projection-identity}.

Apply Eq.~\eqref{eq:projection-identity} to the block $\{1,2,3\}$, written
as the orthonormal basis $\{v_3,v_2,v_1\}$.  Since $v_4\perp v_3$ by the
block $\{3,4,5\}$ and $v_{12}\perp v_1$ by the block $\{11,12,1\}$, we get
\begin{equation}\label{eq:g32-id-A}
    \langle v_4,v_{12}\rangle
    =\langle v_4,v_2\rangle\langle v_2,v_{12}\rangle.
\end{equation}
With the parametrization above,
\[
    \langle v_4,v_{12}\rangle=\overline b\,\overline c,
    \qquad
    \langle v_4,v_2\rangle=\overline a\,q,
    \qquad
    \langle v_2,v_{12}\rangle=-\overline p\,\overline d.
\]
Thus
\[
    \overline b\,\overline c
    =-\overline a\,\overline d\,q\,\overline p,
\]
and after complex conjugation,
\begin{equation}\label{eq:g32-constraint-1}
    bc=-ad\,\overline q\,p.
\end{equation}

Next apply Eq.~\eqref{eq:projection-identity} to the block $\{7,8,9\}$,
written as the orthonormal basis $\{v_9,v_8,v_7\}$.  Since
$v_{10}\perp v_9$ by the block $\{9,10,11\}$ and $v_6\perp v_7$ by the
block $\{5,6,7\}$, we obtain
\begin{equation}\label{eq:g32-id-B}
    \langle v_{10},v_6\rangle
    =\langle v_{10},v_8\rangle\langle v_8,v_6\rangle.
\end{equation}
The explicit inner products are
\[
    \langle v_{10},v_6\rangle=ad,
    \qquad
    \langle v_{10},v_8\rangle=-b\,\overline p,
    \qquad
    \langle v_8,v_6\rangle=-qc.
\]
In the last equality the first component of $v_8$ is conjugated once more by
$\langle x,y\rangle=x^\dagger y$; hence the factor is $q$, not $\overline q$.
Therefore
\begin{equation}\label{eq:g32-constraint-2}
    ad=bc\,\overline p\,q.
\end{equation}

Substituting Eq.~\eqref{eq:g32-constraint-1} into
Eq.~\eqref{eq:g32-constraint-2} gives
\[
    ad=(-ad\,\overline q\,p)\overline p\,q
      =-ad\,|p|^2|q|^2.
\]
Hence
\[
    ad\bigl(1+|p|^2|q|^2\bigr)=0.
\]
The factor $1+|p|^2|q|^2$ is strictly positive, so
\[
    ad=0.
\]

If $a=0$, then $|b|=1$ and
\[
    v_4=\begin{pmatrix}0\\0\\b\end{pmatrix}=b e_3.
\]
Thus $[v_4]=[v_{15}]$, contradicting faithfulness.  If $d=0$, then
$|c|=1$ and
\[
    v_6=\begin{pmatrix}c\\0\\0\end{pmatrix}=c e_1.
\]
Thus $[v_6]=[v_{13}]$, again contradicting faithfulness.  Therefore no
faithful orthogonal representation of $G_{32}$ in $\C^3$ exists.
\end{proof}

\begin{remark}
The proof never uses the chromatic number of $G_{32}$.  The contradiction
comes from the propagation of transition amplitudes through the incidence
geometry.  The two factorizations forced by opposite parts of the outer
cycle are algebraically incompatible with distinctness of the central-spoke
rays.  Thus the obstruction is geometric rather than chromatic.
\end{remark}

\section{Conclusion}

The existence of a faithful orthogonal representation, the existence of a
separating set of two-valued states, and the existence of a globally
consistent spectral coloring are distinct questions.  The strong chromatic
condition $\chi(H)>n$ rules out chromatic completeness in dimension $n$,
because it rules out a global assignment of the same $n$ eigenvalue labels
to every context.  It does not rule out a faithful orthogonal representation.

The completed Yu--Oh hypergraph makes this explicit: it has
$\chi=4>3$ but is edge-faithfully represented by 25 rays in $\R^3$.
Greechie's $G_{32}$ has the same strong chromatic number and a separating and
unital set of two-valued states, but no faithful representation in $\C^3$.
Thus even the common value $\chi=4$ does not determine coordinatizability.
The obstruction in $G_{32}$ is detected by direct geometric analysis, not by
coloring or by the mere presence or absence of two-valued states.

Thus chromatic contextuality concerns global consistency of spectral labels,
whereas geometric coordinatizability concerns realizability of incidence
relations by rays.  These two layers of structure should not be conflated.


\begin{acknowledgments}
I thank Mohammad Hadi Shekarriz for valuable comments and suggestions.

The author used OpenAI Codex (GPT-5) for editorial and structural revision of
the manuscript.  The author directed this use, independently verified the
mathematical claims and references, and accepts full responsibility for the
final content.

This research was funded in whole or in part by the \textit{Austrian Science Fund (FWF)} [Grant \href{https://doi.org/10.55776/PIN5424624}{digital object identifier (DOI): 10.55776/PIN5424624}].
The author acknowledges TU Wien Bibliothek for financial support through its Open Access Funding Programme.
\end{acknowledgments}

\bibliography{svozil}

\end{document}
